%% author - Matúš Fignár - xfignam00

\documentclass[a4paper, 11pt]{article}
\usepackage[left=2cm, top=3cm, text={17cm,24cm}]{geometry}
\usepackage[utf8]{inputenc}
\usepackage[T1]{fontenc}
\usepackage{times}
\usepackage[english,czech,slovak]{babel}
\usepackage{csquotes}
\usepackage[backend=biber,style=iso-numeric]{biblatex}
\usepackage{hyperref}

\addbibresource{odkazy.bib}
\begin{document}

\begin{titlepage}
    \begin{center}
        \textsc{\Huge{Vysoké učení technické v~Brně\\[0.4em]}
        \huge{Fakulta informačních technologií}}\\        
        \vspace{\stretch{0.382}}
        \LARGE{Typografie a~publikování\,--\,4. projekt}\\
        \Huge{Bibliografické citácie}\\ %možno zmeniť na "Bibliografické citace"
        \vspace{\stretch{0.618}}
    \end{center}
    \Large{\today\hfill Matuš Fignár}
\end{titlepage}

\begin{center}
    \Huge{\textbf{Psychológia a~fonty}}
\end{center}

\section{Úvod}
Pri čítaní kníh, časopisov alebo článkov existuje mnoho faktorov, ktoré ovplyvňujú čitateľský zážitok. Mnoho ľudí by 
povedalo, že hlavnú úlohu zohrávajú obrázky alebo rozloženie textu, no často sa zabúda na význam samotného štýlu písma. 
Typografia nie je len o~čitateľnosti a~estetickom rozvrhnutí textu, ale aj o~tom, ako jednotlivé fonty ovplyvňujú 
naše emócie a~rozhodovanie.\cite{lupton}

\section{Dôveryhodnosť a~emócie fontov}
Výber fontu môže mať zásadný dopad na dôveryhodnosť textu – 
formálny dokument napísaný písmom Comic Sans čitateľa skôr 
odradí, zatiaľ čo profesionálne písma ako Times New Roman a~Garamond môžu zvýšiť dôveru v~obsah.\cite{designmodo}
Pri výbere fontu by sme sa mali zamyslieť, na akú skupinu ľudí mierime a~čo chceme povedať. 
Písmo nie je len prostriedkom komunikácie, ale aj nástrojom, 
ktorým autor ovplyvňuje náladu a~správanie čitateľa.
Rôzne fonty evokujú rôzne pocity – napríklad serifové 
písma ako Times New Roman sú vnímané ako konzervatívne, zatiaľ čo sans-serif písma ako Helvetica pôsobia 
modernejšie.\cite{garfield} 

\section{Štúdia}
Psychologické účinky písma boli overované aj empiricky. Štúdia uverejnená v~British Journal of Psychology
preukázala, že čitateľská rýchlosť a~porozumenie sa môžu líšiť v~závislosti od použitého písma.\cite{BPS}
Ďalšia štúdia poukázala na to, že naše vnímanie písma je špecifické – nejde len o~tvar písmen, ale o~ich súvislosti s~predchádzajúcimi 
skúsenosťami a~vizuálnym učením.\cite{dyson} Zaujímavým poznatkom je, že menej pohodlný text na čítanie zvyšuje pozornosť a~čitateľ sa
sústredí viac na text.\cite{oppenheimer}

\section{Font a~myšlienka textu}
Čitateľ má tendenciu spájať rôzne typy písma so zámerom textu alebo osobnostnými črtami autora – napríklad písmo typu Brush Script
je vnímané ako ženské a~emocionálne, zatiaľ čo Arial ako neutrálne.\cite{hlavnickova}
Podstatné je však podotknúť, že hoci fonty majú emocionálnu a~psychologickú väzbu 
s~ich autorom, nedá sa na ich základe vyvodiť plnohodnotný profil autora.\cite{Fajkusova}

\section{Fonty v~marketingu}
Fonty majú svoj význam aj v~marketingovom prostredí. Výrobok s~dobre zvoleným fontom môže mať 
v~priemere až 2-krát väčšiu úspešnosť a~v~niektorých prípadoch až 4-krát väčšiu ako výrobok s~nehodiacim sa fontom.\cite{blueghost}
Font je pre slová ako oblečenie pre človeka. Tak ako si človek vyberá oblečenie podľa sociálnej skupiny, v~ktorej sa nachádza,
tak isto by mala značka zvoliť, aký font použije.\cite{Grigerova} Font je podstatný pre celkový imidž, ktorý značka chce vytvoriť – napríklad
ručne písané názvy dávajú zákazníkovi priateľský pocit.\cite{Izadi}

\section{Záver}
Dôležitosť výberu písma sa teda netýka len dizajnérov alebo typografov. Je to relevantná otázka pre marketingových 
odborníkov, UX dizajnérov či autorov odborných textov. Výber fontu nie je len estetické rozhodnutie – je 
to psychologický nástroj, ktorý ovplyvňuje vnímanie, porozumenie a~dôveru čitateľa.


\newpage
\printbibliography

\end{document}
